\documentclass[a4paper,12pt]{article}

\usepackage[a4paper, margin=2.5cm]{geometry}
\usepackage{graphicx}
\usepackage[bahasa]{babel}
\usepackage{xcolor}
\usepackage[normalem]{ulem}
\usepackage{amsmath} 
\usepackage{amssymb} 
\usepackage{enumitem} 

\title{MODUL O}
\author{Diga Firlana}
\date{April 2025}

\begin{document}


\begin{center}
    {\LARGE\bfseries LATIHAN OSN-K 2025} \\[1cm]
    {\large Diga Firlana} \\[1cm]
    {\large April 2025} \\[2cm]
    
    \includegraphics[width=0.5\textwidth]{images (3).jpeg}
\end{center}
\newpage

\section*{\color{orange}§1. \color{black}Aljabar}

\begin{enumerate}[label=\arabic*., leftmargin=*, align=left, itemsep=1.2cm]

    \item Diberikan bilangan real $a$ dan $b$ sehingga
    \[
        a + \lfloor b \rfloor = 2{,}16, \qquad 
        b + \lfloor a \rfloor = 2{,}19.
    \]
    Tentukan nilai dari $100(a+b)$.
    
    \item Jika $x$, $y$, $z$ adalah bilangan asli sehingga
    \[
        \frac{2025}{2022} = 1 + \frac{1}{z + \dfrac{x}{y}},
    \]
    tentukan nilai dari $x + y + z$.
    
    \item Buktikan bahwa 
    \[
        x^2 + \frac{8}{xy} + y^2 \geq 8
    \]
    untuk $x,y$ bilangan real positif.
    
    \item Diberikan $a, b, c > 0$ memenuhi $a + b + c = 1$.  
    Tentukan nilai maksimum dari
    \[
        \sqrt{2025a+1} + \sqrt{2025b+1} + \sqrt{2025c+1}.
    \]
    
    \item \textbf{(KTOM Maret 2025)} Tentukan banyaknya pasangan bilangan real $(a, b)$ yang memenuhi
    \[
        a^3 - b = 2a, \qquad 
        b^3 - a = 2b.
    \]
    
    \item \textbf{(KTOM Maret 2025)} Diberikan bilangan real berbeda $a, b, c$ yang memenuhi  
    \[
        36(a - b)^2 = 9(b - c)^2 = 4(c - a)^2.
    \]
    Tentukan nilai dari
    \[
        \left| 36 \left( 
        \frac{a - b}{b - c} + 
        \frac{b - c}{c - a} + 
        \frac{c - a}{a - b} \right) \right|.
    \]
    
    \item \textbf{(TO OSN-K Ruang Belajar)}  
    Diberikan $x, y > 0$ yang memenuhi $(x + y)^2 = 12xy$.  
    Jika
    \[
        \frac{x^4 + y^4}{(x + y)^4} = \frac{m}{n},
    \]
    dengan $m,n$ relatif prima, tentukan nilai $m + n$.
    
    \item \textbf{(KTOM September 2024)}  
    Jika $x$ memenuhi
    \[
        \frac{2}{x} + \frac{1}{x - 1} = 1,
    \]
    maka nilai dari $\dfrac{2}{x} + x$ adalah ...
    
    \item \textbf{(OSN-K 2024)}  
    Diketahui $a, b, c > 0$ memenuhi
    \[
        a + b + c = \frac{32}{a} + \frac{32}{b} + \frac{32}{c} = 24.
    \]
    Nilai terbesar dari $\dfrac{a^{2} + 32}{a}$ adalah ...
    
    \item \textbf{(OSN-K 2024)}  
    Misalkan $x, y > 0$ dengan $x > y$ dan
    \[
        x^2 + y^2 = \left( \frac{545}{272} \right) xy.
    \]
    Tentukan nilai dari $\dfrac{x+y}{x-y}$.

\end{enumerate}

\section*{\color{orange}§2. \color{black}Teori Bilangan}

\begin{enumerate}[label=\arabic*., leftmargin=*, align=left, itemsep=1.2cm, resume]

    \item Tentukan nilai dari $\sum\limits_{n=1}^{2025} \operatorname{FPB}(n,3)$.

    \item Tentukan banyak bilangan 4-digit yang habis dibagi 2 tetapi tidak habis dibagi 6.

    \item \textbf{(TO Ambis Camp)}  
    Misalkan $a$ dan $b$ bilangan asli memenuhi
    \[
        3a + 7b = 2025.
    \]

    \item \textbf{(KTOM Maret 2025)}  
    Tentukan jumlah seluruh bilangan asli $n$ yang merupakan kelipatan 15 dan memiliki tepat 15 faktor.

    \item Digit kedua dari $2025^{2025}$ adalah ...

    \item \textbf{(KTOM Maret 2025)}  
    Tentukan nilai dari
    \[
        \operatorname{FPB}(20^3 + 50, 20^5 + 50, 20^7 + 50, \ldots, 20^{2025} + 50).
    \]

    \item Tentukan hasil dari
    \[
        11^2 - 1^2 + 12^2 - 2^2 + 13^2 - 3^2 + \cdots + 20^2 - 10^2.
    \]

    \item Tentukan sisa pembagian dari $1^3+2^3+3^3+\cdots+2025^3$ oleh $1000$.

    \item Tentukan sisa pembagian dari $64^{2025} + 20^{2024}$ oleh $9$.

    \item Banyaknya bilangan yang habis dibagi 3 dan tidak memuat angka 6 adalah ...
    
\end{enumerate}

\end{document}

